\chapter{困ったときは}\label{chap:trouble}

\newcommand{\qa}[1]{\subsection*{\textcolor{lsred}{Q.}\ #1}\addcontentsline{toc}{subsection}{#1}}

% ------------------------------------------------------------------------------
\section{警告の帯が出た}

注意すべきことがあると、プレビューの上に帯が出ます。帯は \ui{×} で閉じられます。

\subsection*{バイトオーダーの注意}\label{sec:byteorder}

\screen[0.62\linewidth]{sun_banner}{バイトオーダーの注意（太陽 Hα の 16bit SER）}{\mk[inw]{banner}{1}\mk[s]{buttons}{2}}

\begin{marklist}
  \item 16bit の SER で、ファイルに記録された数値の並び（バイトオーダー）と、\LS が中身から判定した並びが違うときに出ます。
        灰色の控えめな帯のときは、\textbf{画像が正常に見えていれば問題ありません}（自動判定のほうを使っています）。
  \item 画像が砂嵐のように壊れて見えるときは、\ui{little}\ui{big} を押して読み直します。品質評価タブの\uil{バイトオーダー}と同じです。
\end{marklist}

\subsection*{そのほかの帯}

\begin{tabularx}{\linewidth}{@{}>{\raggedright\arraybackslash}p{0.36\linewidth} X@{}}
\toprule
帯の内容 & どうするか\\
\midrule
\uil{12bitデータが16bitコンテナに入っている可能性があります} & 画像がとても暗いなら \ui{12bit} を押します（品質評価タブの\uil{ビット深度の解釈}と同じ）。\\
\uil{N 枚を自動除外しました（視野外・追跡失敗）} & \ui{詳細…} で、除外したフレームと理由の一覧を見られます。数枚なら気にしなくてかまいません。
  多すぎるときは第\ref{chap:quality}章の「詳細設定（追跡失敗の判定）」を見直します。\\
\uil{MOV・MP4 は macOS のデコーダで読みます…} & 知らせだけです（第\ref{sec:movie}節）。\\
\bottomrule
\end{tabularx}

% ------------------------------------------------------------------------------
\section{よくある質問}

\qa{スタックした結果がぼやける}
\begin{itemize}[itemsep=1pt]
  \item \textbf{仕上げのウェーブレット}を掛けましたか？ スタックの結果は、ふつう少しぼやけて見えます。ウェーブレットで細部を引き出すのが前提です（第\ref{sec:wavelet}節）。
  \item スタックタブの採用率を\textbf{下げて}（例 10\% → 5\%）\ui{再スタック} してみてください。
  \item 位置合わせ領域の大きさを1段小さくして \ui{再アライメント} してみてください。
  \item ピントや撮影時のゆらぎそのものが原因のこともあります。品質グラフの\uil{時系列}で、よい区間だけに\uil{使うフレームの範囲}を絞るのも手です。
\end{itemize}

\qa{ざらざら（ノイズ）が多い}
\begin{itemize}[itemsep=1pt]
  \item スタックタブの採用率を\textbf{上げて}、重ねる枚数を増やします（第\ref{sec:setting-aptop}節）。
  \item ウェーブレットの強調を下げるか、Layer 1〜2 の\uil{ノイズ}を上げます。
  \item 撮影では、露出（ゲイン）を少し上げる、フレーム数を増やすのが効きます。
\end{itemize}

\qa{画像の一部がゆがむ・二重になる・崩れる}
\begin{itemize}[itemsep=1pt]
  \item アライメントタブの詳細設定で\uil{一致度の下限}を上げて（0.6〜0.7）、合わせ間違いを使わないようにします。
  \item 位置合わせ領域の大きさを1段大きくします。
  \item 位置ずれの探索範囲が狭すぎないか確かめます（像が大きく踊る日は ±24〜±32）。
\end{itemize}

\qa{惑星の縁や暗い所だけぼやける・位置合わせ領域が置かれない}
アライメントタブの詳細設定で\uil{配置する模様の強さ}・\uil{配置する明るさ}を下げるか、\uil{配置を編集}で手で足します（第\ref{sec:apedit}節）。

\qa{色がおかしい（緑っぽい・紫っぽい・格子模様が出る）}
\begin{itemize}[itemsep=1pt]
  \item カメラの RAW は緑がかって読まれます。仕上げの \ui{自動（灰色仮説）} を押してください。
  \item 格子模様や、ありえない色になるときは、品質評価タブの\uil{色の配列（Bayer）}が間違っている可能性があります。4種類を試してください。
  \item 惑星の縁の赤・青のにじみは、RGB チャンネル合わせの \ui{自動で合わせる} で取れます。
\end{itemize}

\qa{惑星の縁の外に黒い輪ができる}
ウェーブレットの強調のしすぎです。Layer 1〜2 の強調を下げるか、\uil{輪抑制}を 0.2〜0.5 に上げます。

\qa{処理が遅い・Mac が重くなる}
\begin{itemize}[itemsep=1pt]
  \item 大きな画像で対象が一部だけなら、\textbf{処理範囲}を使います（第\ref{sec:roi-howto}節）。
  \item 位置ずれの探索範囲を狭くする（±8〜±12）、\uil{参照の反復精密化}をオフにすると、アライメントが速くなります。
  \item メモリが少ない Mac では、スタックタブの\uil{低メモリモード}をオンにします。ドリズルの倍率を上げるほどメモリを使います。
  \item 採用率だけを変えるなら \ui{再スタック} だけで済みます。品質評価からやり直す必要はありません。
\end{itemize}

\qa{ドリズルの診断で、いつも 1× と言われる}
焦点距離が十分に長く、すでに細かく写っている撮影では、ドリズルの効果はほとんどありません。これは正常です。
診断の理由の行で\uil{（いちばんの制約）}が付いたものを見てください。\uil{採用枚数}が制約なら、フレーム数を増やすか採用率を上げると変わることがあります。

\qa{設定が消えた}
\LS は起動するたびに設定を初期値に戻します。よく使う設定はプリセットに保存してください（第\ref{sec:preset}節）。
解析の結果（サイドカー）は残っているので、同じファイルを開けばスタックから始められます。

\qa{書き出した画像が、画面より暗い（明るい）}
\ui{表示} メニューの \uil{表示を明るくする} がオンだと、画面の上でだけ明るく見せています。オフにすると、書き出すファイルと同じ明るさになります。
明るさはレベル補正（第\ref{sec:levels}節）で調整してから書き出してください。

\qa{［アライメント］や［スタック］が押せない}
前の工程が終わっていないか、前の工程に関わる設定を変えた直後です。状態表示に「品質評価からやり直してください」などと出ているので、その工程から押し直してください。
